%% ma: L10-B12-0020
%% lop: 10
%% chuong: 5
%% bai: 12
%% dang: 10.12.7
%% loai: TL
%% muc: VD
%% dap_an: "-"
%% nguon: "chuyen-de-cac-so-dac-trung-cua-mau-so-lieu-khong-ghep-nhom-toan-10-chuong-trinh-moi.pdf (GV. Phan Nhật Linh) (P01), Bài 01 Bài tập tự luận 20"
%% kien_thuc: "Lớp 10 Bài 12 (số gần đúng, sai số tuyệt đối, độ chính xác, sai số tương đối, quy tròn số)"
%% trang_thai: chinh-thuc
%% ghi_chu: "Bỏ tên thương hiệu và mã thép trong đề gốc (không ảnh hưởng dữ kiện). Tự duyệt theo yêu cầu giáo viên 22/10/2026."
\begin{ex}
Nhà sản xuất công bố chiều dài và chiều rộng của một tấm thép hình chữ nhật lần lượt là $15\pm0{,}05$ cm và $6\pm0{,}05$ cm. Hãy tính diện tích tấm thép và đánh giá sai số tuyệt đối của diện tích đó.
\loigiai{Gọi $\overline a,\overline b$ là chiều dài và chiều rộng thực: $14{,}95\le\overline a\le15{,}05$; $5{,}95\le\overline b\le6{,}05$.

Suy ra $88{,}9525\le\overline S=\overline a\,\overline b\le91{,}0525$.

Diện tích gần đúng $S=15\cdot6=90$, nên $-1{,}0475\le\overline S-90\le1{,}0525$, tức là $\left|\overline S-90\right|\le1{,}0525$.

Vậy diện tích tấm thép là $90\pm1{,}0525$ (cm$^2$).}
\end{ex}
